\documentclass[10pt,a4paper]{amsart}
\usepackage[margin=19mm]{geometry}
\usepackage{amsmath,amssymb,mathtools}
\usepackage[scr=rsfs,cal=euler]{mathalfa}
\usepackage{xfrac}
\usepackage[unicode,hidelinks,bookmarks=false]{hyperref}
\newcommand{\ie}{i.e.\ }
\newcommand{\cf}{cf.\ }
\newcommand{\resp}{respectively\ }
\newcommand{\Subsection}{\subsection}
\providecommand{\nobreakdash}{\nobreak\hbox{-}}
\setlength{\emergencystretch}{2em}
\multlinegap0pt
\usepackage{amssymb}
\usepackage{xcolor}
\newtheorem{proposition}{Proposition}
\title{Hilbert space models and Blaschke frames}
\author{Connor Evans}
\date{Source: arXiv:2510.18816v3 (CC BY 4.0)}
\begin{document}
\maketitle

\noindent\emph{Source and licence.} Hilbert space models and Blaschke frames. Version arXiv:2510.18816v3 (CC BY 4.0), distributed by arXiv under the Creative Commons Attribution 4.0 International licence, http://creativecommons.org/licenses/by/4.0/, as recorded in the arXiv licence metadata for this exact version and shown on the abstract page https://arxiv.org/abs/2510.18816v3. Full licence notice: \texttt{licenses/arxiv-2510.18816-CC-BY-4.0.txt}.\par\smallskip

\noindent\textit{Editorial scope.} Counted unit: the article's proposition Psurjective (source0.tex lines 384--391). The article's complete original proof of it is delivered with it (source0.tex lines 393--437, one \texttt{proof} environment of 45 lines). No mathematics is written, repaired or re-derived: the statement and the proof are the article's own lines, in the article's own notation, and no retained line is substituted or re-flowed.  Retained uncounted as the source-local context the proof needs: lines 252--254.  Nothing was omitted from any retained span, so the package carries no omission record.  No retained line contains a citation, so the wrapper carries no bibliography and no bibliography entry.  No retained line contains a figure environment, an image-inclusion command or drawing code, so no image is delivered and the package declares no asset dependency.  Editorial formatting only, in addition: an AMS \texttt{amsart} preamble that restates the article's own declarations and macros used by the retained lines, namely the environments \texttt{proposition} on separate counters, together with the standard packages those lines need.  The retained environments print in this document as Proposition 1 in order of appearance, whereas the article prints the same items with the numbering of its own full document.  The complete native source of the article as distributed in the arXiv e-print for this exact version is bundled unchanged as \texttt{sources/187529/source0.tex}.
\begin{equation}\label{factorisation}
    p(z) = \prod_{j=1}^{d}(z+\lambda_{j})~~~\text{and}~~~\overline{q}(z)=\prod_{j=1}^{d}(1+\overline{\lambda_{j}}z)
\end{equation}

\begin{proposition}\label{Psurjective}
Let $p(z)=z^{d}+c_{1}z^{d-1}+\hdots+c_{d}$ for all $z\in\overline{\mathbb{D}}$ with coefficients $(c_{k})_{k=1}^{d}\in\mathbb{G}_{d}$. For an isometry $V$ on a Hilbert space $\mathcal{H}$, define the operator $p(V^{*})$ on $\mathcal{H}$ by

$$p(V^{*})=(V^{*})^{d}+c_{1}(V^{*})^{d-1}+\hdots +c_{d}.$$

\noindent Then $p(V^{*})$ is surjective.

\end{proposition}

\begin{proof}
The polynomial $p$ has the factorisation given by equation (\ref{factorisation}). Hence, there exists scalars $\lambda_{1},\hdots,\lambda_{d}\in\mathbb{D}$ such that

\begin{equation}\label{inequalityforinv}
0\leq \dfrac{2\lvert\lambda_{j}\rvert}{1+\lvert\lambda_{j}\rvert^{2}}<1
\end{equation}

\noindent for $j=1,\hdots,d$. As $V$ is an isometry on $\mathcal{H}$, we can write the operator $(V^{*}+\lambda_{j})(V+\overline{\lambda}_{j})$ as

$$
    (V^{*}+\lambda_{j})(V+\overline{\lambda}_{j}) = 1_{\mathcal{H}} + 2\mathrm{Re}(\lambda_{j} V) + \lvert \lambda_{j} \rvert^{2} =(1_{\mathcal{H}}+\lvert\lambda_{j}\rvert^{2})\bigg(1_{\mathcal{H}}+\dfrac{2\mathrm{Re}(\lambda_{j} V)}{1+\lvert\lambda_{j}\rvert^{2}}\bigg).$$

\noindent Since

$$ || 2\mathrm{Re}(\lambda_{j} V) ||_{\mathcal{B}(\mathcal{H})} = || \lambda_{j} V + \overline{\lambda_{j}} V^{*} ||_{\mathcal{B}(\mathcal{H})} \leq 2 \lvert\lambda_{j} \rvert, $$

\noindent we have that by inequality (\ref{inequalityforinv}), 

\begin{equation}\label{inequalityforinv2}
\Bigg\|\dfrac{2\mathrm{Re}(\lambda_{j} V)}{1+\lvert\lambda_{j}\rvert^{2}}\Bigg\|_{\mathcal{B}(\mathcal{H})}\leq \dfrac{2\lvert\lambda_{j}\rvert}{1+\lvert\lambda_{j}\rvert^{2}} < 1,
\end{equation}

\noindent hence the operator $(V^{*}+\lambda_{j})(V+\overline{\lambda}_{j})$ is invertible, as each of the operators $1_{\mathcal{H}}+\lvert\lambda_{j}\rvert^{2}$ and

$$ 1_{\mathcal{H}}+\dfrac{2\mathrm{Re}(\lambda_{j} V)}{1+\lvert\lambda_{j}\rvert^{2}}$$ 

\noindent are invertible for all $\lambda_{j}\in\mathbb{D}$ and $j=1,\hdots,d$.

 A bounded operator is surjective if and only if it has a bounded right inverse. For $\lambda_{j}\in\mathbb{D}$ for $j=1,\hdots,d$, we claim that the operator $V^{*}+\lambda_{j}$ has a bounded right inverse, denoted by $\dagger$, given by


$$(V^{*}+\lambda_{j})^{\dagger}=(V+\overline{\lambda}_{j})\bigg[(V^{*}+\lambda_{j})(V+\overline{\lambda}_{j})\bigg]^{-1}.$$

\noindent Indeed, we have proven invertibility of the square bracket for the operator given above and boundedness can be proven by a Neumann series argument and an application of the triangle inequality. By direct calculation,


$$(V^{*}+\lambda_{j})(V^{*}+\lambda_{j})^{\dagger}
    =(V^{*}+\lambda_{j})(V+\overline{\lambda}_{j})\bigg[(V^{*}+\lambda_{j})(V+\overline{\lambda}_{j})\bigg]^{-1} 
    = 1_{\mathcal{H}}.$$

\noindent Therefore the operator $V^{*}+\lambda_{j}$ is surjective. Furthermore, the operators

$$p(V^{*}) = \prod_{j=1}^{d}(V^{*}+\lambda_{j})$$

\noindent are surjective. This comes as consequence of the composition of surjective operators being surjective.\end{proof}

\end{document}
